\section{Анализ предметной области и формирование требований к гибридной интеллектуальной системе}

\subsection{Характеристика предметной области}

Предметной областью разрабатываемой системы является взаимодействие пользователя
с интерактивной виртуальной средой посредством естественного языка. В качестве
конкретной среды выбрана Minecraft версии 1.20.1 — песочница с открытым миром,
широко применяемая в исследованиях интеллектуальных агентов в качестве
испытательного полигона~\cite{Wang2024,Fan2022}. Выбор обусловлен тремя обстоятельствами: среда
обладает формализованными и доступными для машинной обработки данными о своих
объектах, допускает программное расширение, а поведение пользователя в ней не
предопределено сценарием, что обеспечивает разнообразие запросов.

Основными сущностями предметной области являются:

\begin{itemize}
  \item \textit{предметы} — объекты, которые могут находиться в инвентаре
        пользователя и расходоваться при изготовлении;
  \item \textit{блоки} — объекты, размещаемые в мире и добываемые из него;
  \item \textit{существа} — подвижные объекты мира, из которых при победе над
        ними извлекаются предметы;
  \item \textit{рецепты} — способы получения одних предметов из других с
        указанием состава, количества и требуемого оборудования;
  \item \textit{группы} — именованные множества взаимозаменяемых предметов,
        позволяющие задавать состав рецепта обобщённо;
  \item \textit{инструменты} — предметы, наличие которых необходимо для добычи
        определённых блоков, характеризуемые видом и уровнем материала.
\end{itemize}

Между сущностями определены отношения: предмет входит в состав рецепта, рецепт
производит предмет, блок даёт предмет при добыче, блок требует инструмента
определённого уровня, группа содержит предметы. Совокупность этих отношений
образует ориентированный граф способов получения объектов, обход которого и
составляет основу решения задач в разрабатываемой системе.

Существенной особенностью предметной области является наличие
\textit{состояния пользователя}: содержимого инвентаря, имеющихся инструментов,
показателей здоровья и параметров окружения. Одна и та же формулировка запроса
при разном состоянии требует разного ответа, поэтому знания о мире и сведения о
текущем состоянии должны использоваться совместно.

Источником знаний служат штатные файлы данных среды, описывающие рецепты,
группы, таблицы добычи и наименования объектов, а также официальная локализация.
Эти данные доступны в машиночитаемом виде и позволяют построить базу знаний
автоматически, без ручного наполнения. Часть сведений, однако, в файлах данных
отсутствует и задаётся программным кодом среды; соответствующие ограничения
рассматриваются в подразделе~1.5.

Далее в пояснительной записке используется следующая терминология.
\textit{Гибридная интеллектуальная система} — система, в которой логико-символьные
и нейросетевые компоненты объединены в единой модели решения задач.
\textit{Решатель задач} — компонент, выполняющий вывод на знаниях базы и
формирующий решение. \textit{План} — упорядоченная совокупность типизированных
шагов, приводящая от текущего состояния пользователя к поставленной цели.
\textit{Вето} — отклонение решателем действия, предложенного языковой моделью,
с указанием причины.

\subsection{Анализ решаемого класса задач}

Разрабатываемая система решает задачи трёх классов.

\textbf{Планирование.} По заданной цели и текущему состоянию пользователя
требуется построить последовательность действий, приводящую к достижению цели.
Условиями задачи являются целевой объект, требуемое количество и содержимое
инвентаря; результатом — план с типизированными шагами либо обоснованное
заключение о недостижимости цели. Задача относится к классической постановке планирования в пространстве
состояний~\cite{Fikes1971,Ghallab2004}. Критерием корректности служит выполнимость
каждого шага плана при условии выполнения предшествующих, критерием качества —
трудоёмкость плана для пользователя.

\textbf{Проверка допустимости.} По предложенному действию требуется установить,
разрешено ли оно правилами предметной области, и при отказе назвать причину.
Условиями задачи являются вид действия, его объект и количество; результатом —
вердикт с объяснением. Критерием корректности является отсутствие ложно
разрешённых действий; отказ в действии, которое могло быть разрешено, считается
менее тяжёлым нарушением, чем выполнение недопустимого.

\textbf{Понимание запроса.} По свободной реплике пользователя требуется определить
намерение и сопоставить упомянутые объекты с каноническими идентификаторами
предметной области. Результатом является структурированное представление запроса
либо уточняющий вопрос, если запрос допускает несколько истолкований.

Совместное применение двух подходов необходимо потому, что ни один из них не
решает всех трёх задач~\cite{Russell2021}. Задача понимания запроса не поддаётся формальному
описанию: множество способов выразить одно и то же намерение на естественном
языке не перечислимо, а разговорные и жаргонные наименования объектов не
содержатся ни в каких справочниках среды. Напротив, задачи планирования и
проверки допустимости требуют гарантированной корректности, которую вероятностная
модель обеспечить не может: она не располагает состоянием инвентаря, не выполняет
вывода по графу способов получения и не проверяет своих утверждений~\cite{Ji2023,Jurafsky2024}.

Таким образом, распределение ролей определяется не соображениями удобства, а
природой задач: нейросетевой компонент отвечает за понимание и формулирование,
логико-символьный — за знание, вывод и гарантии.

\subsection{Анализ базы знаний, решателя задач и интерфейса}

В соответствии с многосеместровой организацией курсового проектирования гибридная
система строится на компонентах, разработанных на предыдущих этапах. Ниже
приведён анализ их фактического состояния и пригодности к включению в систему.

\textbf{Интерфейс (6-й семестр).} Реализован в виде модификации игрового клиента
на платформе Fabric (язык Java, версия среды 1.20.1) и серверной части на базе
фреймворка FastAPI. Модификация обеспечивает приём реплик пользователя из игрового
чата, отображение ответов, внутриигровой экран-оверлей и передачу телеметрии:
событий состояния, фактов победы над существами и снимков инвентаря. Серверная
часть реализует цикл вызова инструментов, при котором языковая модель выбирает
имя действия и передаёт запрос в свободной форме, а исполняемая команда собирается
сервером. Компонент работоспособен, покрывает требуемые функции и включается в
гибридную систему \textbf{без принципиальных изменений}; доработке подлежит лишь
вывод плана и объяснений в оверлей.

\textbf{База знаний (4-й семестр).} Представлена справочником объектов в формате
JSON, содержащим 261~запись (211~уникальных идентификаторов при 50~дублирующихся)
с русскими и английскими наименованиями. Записи справочника \textbf{не связаны
между собой отношениями}: он допускает лишь сопоставление наименования с
идентификатором. Рецепты хранятся в отдельном файле, ключами и значениями ячеек
в котором выступают русские наименования, не связанные с идентификаторами
объектов; выходное количество не указано. Вывод по таким данным невозможен, они
пригодны лишь для отображения пользователю. Требование к базе знаний,
предъявляемое на 4-м семестре, — наличие декларативных знаний, задающих сущности
и связи между ними, — не выполняется. Компонент подлежит \textbf{разработке
заново}; накопленные вручную разговорные наименования объектов при этом
сохраняются и переносятся в новую базу знаний как дополнительные метки.

\textbf{Решатель задач (5-й семестр).} Отсутствует. Функции, внешне относимые к
символьной обработке, сводятся к нечёткому сопоставлению строк по алгоритму
поиска наибольшей общей подпоследовательности, списку стоп-слов, выделению
количества регулярным выражением и двум константам-ограничителям. Логический
вывод не выполняется. Компонент подлежит \textbf{разработке заново}.

Для получения объективной оценки исходного состояния был сформирован размеченный
набор из 56~реплик пользователя, разбитый на группы типичных, граничных и
отрицательных случаев, и выполнен прогон существующей системы. Результаты
приведены в таблице~\ref{tab:baseline}.

\begin{table}[h]
  \centering
  \caption{Результаты прогона исходной системы на размеченном наборе}
  \label{tab:baseline}
  \begin{tabular}{|p{0.62\textwidth}|p{0.24\textwidth}|}
    \hline
    Показатель & Значение \\ \hline
    Совпадение с разметкой & 53,6~\% \\ \hline
    Доля небезопасных расхождений & 10,7~\% \\ \hline
    Не распознано как запрос действия & 11 случаев \\ \hline
    Подмена объекта при сопоставлении & 4 случая \\ \hline
    Необработанные словоформы & 2 случая \\ \hline
    Время отклика, медиана & 22~мс \\ \hline
  \end{tabular}
\end{table}

Под небезопасными расхождениями понимаются случаи, в которых система уверенно и
без предупреждения выполняет не то действие, которое требовалось. Этот показатель
существеннее общей точности: отказ пользователь видит и может переформулировать
запрос, тогда как подмена объекта внешне неотличима от корректной работы.

Анализ расхождений выявил четыре группы причин.

\textit{Распознавание намерения по списку ключевых слов} не покрывает естественную
речь: формулировки «нужен хлеб», «хочу стейк», «мне бы десять наковален» не
опознаются как запросы вообще. Это ограничение не пробел в списке, а свойство
подхода, поскольку перечислить все способы выражения намерения невозможно.

\textit{Нечёткое сопоставление строк подменяет объект.} Запрос «дай мне пять
алмазов» приводит к выдаче блока алмазов, поскольку словоформа «алмазов» совпадает
с фрагментом наименования «блок алмазов» и перевешивает правильный вариант.

\textit{То же сопоставление срабатывает на несуществующих объектах.} Запрос
«призови дракона Края» приводит к призыву эндермена, поскольку слово «Края»
встречается в наименовании «странник Края». Показательно, что данное существо
отсутствует в справочнике, однако это не защищает от ошибки: отсутствие записи
не равносильно запрету.

\textit{Словоформы обрабатываются вручную.} Наименования, для которых не написано
специального правила усечения, не распознаются.

Дополнительно выявлены два дефекта данных. Восемь идентификаторов зелий содержат
дополнительные данные в фигурных скобках и не проходят проверку формата
идентификатора, вследствие чего соответствующие объекты не могут быть выданы,
причём сообщение об ошибке указывает неверную причину. В справочнике присутствуют
50~дублирующихся идентификаторов с несовпадающими наборами наименований, что
влияет на ранжирование вариантов при сопоставлении.

\subsection{Анализ существующих гибридных интеллектуальных систем}

\textbf{Нейросимвольные архитектуры.} Направление, объединяющее обучаемые модели
с символьным представлением знаний, оформилось как самостоятельное и
рассматривается в качестве третьей волны развития искусственного
интеллекта~\cite{Garcez2023}. Типовым решением является разделение ответственности, при котором
нейросетевая часть выполняет восприятие и выдвижение гипотез, а символьная —
проверку, структурирование и объяснение. Преимуществом является проверяемость
результата, ограничением — необходимость согласования представлений данных на
границе компонентов. Подход применим к рассматриваемой задаче непосредственно и
принят в качестве основного.

\textbf{Архитектура RAG.} Схема дополнения генерации извлечением~\cite{Lewis2020} предполагает
поиск релевантных фрагментов во внешнем хранилище и передачу их языковой модели
в качестве контекста. Достоинством является опора ответа на проверяемые источники,
ограничением — то, что итоговый ответ по-прежнему порождается моделью и не
подлежит формальной проверке. Для задач планирования и проверки допустимости,
где требуется гарантия, схема недостаточна, однако применимые её элементы —
семантический поиск по наименованиям — используются в разрабатываемой системе.

\textbf{Системы вызова инструментов.} Подход, при котором языковая модель выбирает
внешнюю функцию и её аргументы~\cite{Schick2023}, а также схемы чередования рассуждения и
действия~\cite{Yao2023} обеспечивают выполнение моделью операций во внешней системе.
Ограничением подхода является то, что корректность аргументов остаётся на
ответственности модели. Разрабатываемая система использует этот механизм, однако
аргументы подвергаются проверке решателем, а исполняемая команда формируется
сервером.

\textbf{Агенты в виртуальных средах.} Система Voyager~\cite{Wang2024} демонстрирует обучаемого
агента в среде Minecraft, накапливающего библиотеку навыков в виде программного
кода. Платформа MineDojo~\cite{Fan2022} предоставляет среду и набор задач для обучения
агентов, а набор данных MineRL~\cite{Guss2019} — записи действий человека для обучения с
подкреплением. Общим для этих работ является сосредоточенность на автономном
поведении агента, тогда как в настоящей работе решается иная задача — поддержка
человека, формулирующего цели на естественном языке. Существенно также, что
перечисленные системы не ставят задачи гарантированного соблюдения ограничений.

\textbf{Экспертные системы и продукционные правила.} Классический подход к
представлению знаний в виде правил~\cite{Giarratano2006} обеспечивает объяснимость вывода и
проверяемость результата, но требует ручного наполнения базы и не работает со
свободными формулировками. В разрабатываемой системе продукционные правила
применяются для описания ограничений, которые не выводятся из данных среды.

\textbf{Семантические технологии.} Технология OSTIS~\cite{Golenkov2021} и языки описания
онтологий~\cite{Antoniou2012} предоставляют средства формального представления знаний
предметной области и вывода на них. Онтологическое представление принято в
работе в качестве формы хранения базы знаний, поскольку оно обеспечивает явное
задание отношений между сущностями и независимость знаний от программного кода.

Сопоставление рассмотренных подходов приведено в таблице~\ref{tab:approaches}.

\begin{table}[h]
  \centering
  \caption{Сопоставление подходов к построению интеллектуальных ассистентов}
  \label{tab:approaches}
  \begin{tabular}{|p{0.24\textwidth}|p{0.21\textwidth}|p{0.21\textwidth}|p{0.22\textwidth}|}
    \hline
    Подход & Свободная речь & Гарантии корректности & Учёт состояния \\ \hline
    Продукционные правила & нет & да & да \\ \hline
    Языковая модель & да & нет & нет \\ \hline
    RAG & да & частично & нет \\ \hline
    Вызов инструментов & да & нет & частично \\ \hline
    Нейросимвольная схема & да & да & да \\ \hline
  \end{tabular}
\end{table}

Ни один из подходов, кроме нейросимвольного, не обеспечивает одновременно всех
трёх свойств, что и обосновывает выбор гибридной архитектуры.

\subsection{Требования к гибридной интеллектуальной системе}

\subsubsection*{Функциональные требования}

\begin{itemize}
  \item система должна решать задачи планирования достижения цели, проверки
        допустимости действий и понимания свободных запросов пользователя;
  \item система должна распознавать намерение из реплики, сформулированной
        произвольно, включая разговорные наименования объектов и их словоформы;
  \item система должна строить план достижения цели с учётом содержимого
        инвентаря пользователя, раскрывая изготовление, добычу и требования к
        инструменту, в том числе рекурсивно;
  \item система должна проверять предложенное действие по правилам базы знаний
        и отклонять недопустимое с указанием причины;
  \item система должна поддерживать два режима работы, различающиеся составом
        доступных действий;
  \item система должна формировать задания пользователю только после
        подтверждения их выполнимости выводом по базе знаний;
  \item система должна сопровождать решение объяснением, представляющим
        последовательность применённых правил;
  \item при неоднозначном запросе система должна задавать уточняющий вопрос,
        при некорректном — отказывать с объяснением, не выполняя молчаливой
        подстановки значений;
  \item при недостижимости цели система должна указывать шаг, на котором вывод
        оказался невозможен.
\end{itemize}

\subsubsection*{Нефункциональные требования}

\begin{itemize}
  \item \textit{корректность}: доля недопустимых действий должна составлять
        ноль вне зависимости от формулировки запроса;
  \item \textit{полнота}: план должен строиться для объектов, способы получения
        которых представлены в данных среды; границы применимости должны быть
        измерены и документированы;
  \item \textit{устойчивость}: система должна сохранять работоспособность при
        некорректных, неполных и противоречивых запросах, а также при замене
        языковой модели на менее мощную;
  \item \textit{объяснимость}: каждое решение должно сопровождаться
        воспроизводимой цепочкой применённых правил;
  \item \textit{производительность}: время вывода решателя не должно превышать
        200~мс на запрос, суммарное время отклика системы — 2~с;
  \item \textit{модульность}: база знаний, решатель, интерфейс и языковая модель
        должны допускать независимую замену; решатель не должен содержать
        сведений о конкретной среде;
  \item \textit{воспроизводимость}: база знаний должна собираться из исходных
        данных автоматически, результаты тестирования — воспроизводиться при
        фиксированных версиях моделей и параметрах.
\end{itemize}

\subsection{Выводы}

Предметной областью работы является взаимодействие пользователя с интерактивной
виртуальной средой на естественном языке; в качестве среды выбрана Minecraft
версии 1.20.1, обладающая машиночитаемым описанием своих объектов. Выделены
сущности предметной области и отношения между ними, образующие граф способов
получения объектов.

Определены три класса решаемых задач: планирование, проверка допустимости и
понимание запроса. Показано, что необходимость объединения двух подходов
обусловлена природой этих задач: понимание свободной речи не поддаётся
формализации, тогда как планирование и проверка требуют гарантированной
корректности, недостижимой вероятностными методами.

Анализ компонентов предшествующих этапов показал, что интерфейс 6-го семестра
работоспособен и включается в систему практически без изменений, тогда как база
знаний в требуемом смысле отсутствует — имеющийся справочник не содержит связей
между сущностями, — а решатель задач отсутствует полностью. Оба компонента
подлежат разработке заново.

Измерение исходного состояния на размеченном наборе из 56~реплик дало 53,6~\%
совпадения с разметкой при 10,7~\% небезопасных расхождений. Установлены четыре
группы причин, из которых две — распознавание намерения по списку ключевых слов
и нечёткое сопоставление строк — являются ограничениями подхода, а не дефектами
реализации. Выявлено, что отсутствие записи об объекте в справочнике не
равносильно запрету на действие с ним.

Рассмотрены существующие решения: нейросимвольные архитектуры, схема RAG, системы
вызова инструментов, агенты в виртуальных средах, экспертные системы и
семантические технологии. Сопоставление показало, что одновременную поддержку
свободной речи, гарантий корректности и учёта состояния обеспечивает только
нейросимвольная схема, которая и принята в качестве основной.

Сформированы функциональные и нефункциональные требования к системе. Ключевым
из них является требование нулевой доли недопустимых действий вне зависимости
от формулировки запроса, достижимое лишь при вынесении проверок за пределы
языковой модели.
